\documentclass{beamer}

\usetheme{Madrid}
\usepackage[utf8]{inputenc}
\usepackage[brazil]{babel}
\usepackage{graphicx}
\usepackage{amsmath}
\usepackage{hyperref}

\title{IA e Programação na Geografia}
\subtitle{Do Geoprocessamento à Inteligência Artificial: Uma Conversa com Jovens Geocientistas}
\author{Prof. MSc. Gabriel Passos}
\institute{Pontifícia Universidade Católica do Paraná (PUCPR) \\ Universidade Federal de Viçosa (UFV) \\ Centro de Ensino Superior de Maringá (CESUMAR)}
\date{Abril - 2026}

\begin{document}

%------------------------------------------------
\begin{frame}
\titlepage
\end{frame}

%------------------------------------------------
\begin{frame}{Contexto da Apresentação}
\begin{itemize}
    \item Bacharel em Geografia (UFV - Turma 2018 - 2022)
    \item Mestre em Computação Aplicada (UEPG, 2022 - 2025)
    \item Doutorando em Ciência da Computação (PUCPR - 2026 - 2028)
    \item Atuação na interface: Geografia + Programação + IA
    \item Atuação na interface: Medicina + IA (alô, Josué de Castro)
    \item Áreas de interesse: Processamento Digital de Imagens e Engenharia de Dados 
\end{itemize}
\end{frame}

%------------------------------------------------
\begin{frame}{Grupo de Pesquisa: Geomorfologia do Quaternário}
\begin{itemize}
    \item Linha de pesquisa em Geomorfologia: processos de modelagem do relevo, dinâmica de paisagens e evolução do território
    \item Climatologia: variabilidade climática, análise de padrões atmosféricos e impactos ambientais
    \item Integração com geotecnologias e sensoriamento remoto
    \item Uso de dados espaciais e séries temporais
    \item Convergência com IA: modelagem preditiva de fenômenos naturais
\end{itemize}
\end{frame}

%------------------------------------------------
\begin{frame}{Desafios na Geomorfologia Contemporânea}
\begin{itemize}
    \item Alta complexidade dos sistemas naturais e não linearidade dos processos geomorfológicos
    \item Escalas temporais e espaciais múltiplas (do evento extremo à evolução de paisagens)
    \item Limitações de dados em campo e incertezas em medições ambientais
    \item Integração entre modelos físicos e modelos estatísticos
    \item Crescente necessidade de abordagens computacionais e inteligência artificial
\end{itemize}
\end{frame}


%------------------------------------------------
\begin{frame}{Por que falar de IA na Geografia?}
\begin{itemize}
    \item Crescimento exponencial de dados geoespaciais
    \item Sensoriamento remoto e imagens de satélite
    \item Necessidade de automação na análise espacial
    \item Decisão baseada em dados (data-driven geography)
\end{itemize}
\end{frame}

%------------------------------------------------
\begin{frame}{Minha base na Geografia (UFV)}
\begin{itemize}
    \item Atuação em Projetos de Pesquisa/Extensão (LASAC/DEM/UFV) - territorializando a saúde em Viçosa
    \item Disciplinas: EST 103, EAM 330, EAM 450, EAM 451, ENF 314, SOL 480, SOL 481, várias disciplinas do DGE, disciplinas do DPI 
    \item Uso de SIG (Sistemas de Informação Geográfica)
    \item Python e R para automação de análises
    \item Manipulação de dados espaciais
\end{itemize}
\end{frame}

%------------------------------------------------
\begin{frame}{Ferramentas que usei na graduação}
\begin{itemize}
    \item Python: automação e análise espacial
    \item R: estatística aplicada à Geografia
    \item QGIS / ArcGIS
    \item: Google Colab
    \item Bibliotecas: raster, sf, geopandas
\end{itemize}
\end{frame}

%------------------------------------------------
\begin{frame}{Transição para Computação (UEPG)}
\begin{itemize}
    \item Ampliação da base matemática e computacional
    \item Machine Learning aplicado a dados ambientais
    \item Modelagem estatística avançada
    \item Integração entre dados espectrais e variáveis do solo
\end{itemize}
\end{frame}

%------------------------------------------------
\begin{frame}{Agricultura de Precisão}
\begin{itemize}
    \item Monitoramento de variabilidade espacial
    \item Uso de sensores e espectroscopia
    \item Previsão de atributos do solo
    \item Tomada de decisão otimizada no campo
\end{itemize}
\end{frame}

%------------------------------------------------
\begin{frame}{Machine Learning na prática}
\begin{itemize}
    \item Regressão e classificação
    \item Random Forest, SVM, Redes Neurais
    \item Validação cruzada
    \item Métricas de desempenho (RMSE, R², MAE, etc)
\end{itemize}
\end{frame}

%------------------------------------------------
\begin{frame}{Exemplo de fluxo de trabalho}
\begin{itemize}
    \item Coleta de dados
    \item Pré-processamento
    \item Extração de características
    \item Treinamento do modelo
    \item Validação e aplicação espacial
\end{itemize}
\end{frame}

%------------------------------------------------
\begin{frame}{Conexão Geografia + IA}
\begin{itemize}
    \item Geografia quantitativa e computacional (discussão importante: o Paradigma Quantitativo de hoje NÃO é o Quantitativo de Christofoletti)
    \item Modelagem de fenômenos espaciais complexos
    \item Automação da análise geográfica
    \item Escalabilidade de estudos ambientais
\end{itemize}
\end{frame}

%------------------------------------------------
\begin{frame}{Impactos acadêmicos e práticos}
\begin{itemize}
    \item Agricultura mais eficiente e sustentável
    \item Redução de custos de análise
    \item Maior precisão em estudos ambientais
    \item Apoio à tomada de decisão
\end{itemize}
\end{frame}

%------------------------------------------------
\begin{frame}{Desafios}
\begin{itemize}
    \item Qualidade e volume de dados
    \item Interpretação dos modelos de IA
    \item Integração interdisciplinar
    \item Infraestrutura computacional
\end{itemize}
\end{frame}

%------------------------------------------------
\begin{frame}{O futuro da Geografia Computacional}
\begin{itemize}
    \item IA generativa aplicada ao espaço geográfico
    \item Gêmeos digitais do território
    \item Monitoramento em tempo real
    \item Geografia cada vez mais automatizada
\end{itemize}
\end{frame}

%------------------------------------------------
\begin{frame}{Mensagem final}
\begin{itemize}
    \item A Geografia está em transformação (todavia ainda defendo que exista uma "Crise na Geografia")
    \item Programação é uma competência essencial do geógrafo moderno (devemos falar sobre as NOVAS competências profissionais do Geográfo)
    \item IA não substitui o geógrafo, mas sim, amplia sua capacidade
\end{itemize}
\end{frame}

%------------------------------------------------
\begin{frame}{Agradecimentos}
\begin{itemize}
    \item Universidade Federal de Viçosa (UFV)
    \item Laboratório de Geomorfologia do Quarternário (DGE/UFV)
    \item Grupo de Pesquisas Geográficas (GPG/DGE/UFV)
    \item Orientadores e professores da graduação e pós-graduação (PPGEO/UFV)
    \item Colegas de pesquisa e laboratório
    \item Familiares e amigos pelo apoio contínuo
    \item Todos os presentes nesta sessão
\end{itemize}
\end{frame}

%------------------------------------------------
\begin{frame}{Contatos}
\begin{itemize}
    \item E-mail: gabriel@codaqui.dev
    \item GitHub: github.com/iamdatagabs
    \item LinkedIn: linkedin.com/in/passosgabriel
    \item Lattes: lattes.cnpq.br/9928924222649468
    \item ORCID: 0000-0001-5813-7462
\end{itemize}
\end{frame}

%------------------------------------------------
\begin{frame}{Obrigado}
\centering
Perguntas?
\end{frame}

\end{document}
